\chapter{System Analysis}

\section{System Analysis}

This chapter presents a comprehensive analysis of the Distributed Job Scheduler system, covering requirement analysis, feasibility analysis, and object-oriented analysis using UML diagrams.

\subsection{Requirement Analysis}

\subsubsection{Functional Requirements}

The functional requirements describe the specific behaviors and capabilities that the system must provide. They are illustrated using a use case diagram and detailed use case descriptions.

\placeholderfig{Use Case Diagram}{Use case diagram showing actors and use cases for the Distributed Job Scheduler}

\begin{description}
    \item[Actor 1: Client (User)] A user who submits stress-test jobs to the system.
    \item[Actor 2: Worker Node] A machine that receives and executes scheduled jobs.
    \item[Actor 3: Administrator] A user responsible for managing authentication tokens.
    \item[Actor 4: Scheduler] The central system component (also an actor for internal interactions).
\end{description}

\textbf{Use Case Descriptions:}

\begin{enumerate}
    \item \textbf{Register Client:} A user registers as a client using a valid authentication token. The system validates the token, creates a client record, and returns a client ID.
    
    \item \textbf{Submit Job:} An authenticated client submits a job specifying the job type and parameters (e.g., CPU cores, memory size, duration). The system creates a job record with status \texttt{"not started"}.
    
    \item \textbf{Get Job Status:} A client queries the status of a previously submitted job. The system returns the current status (e.g., \texttt{"not started"}, \texttt{"ongoing"}, \texttt{"completed"}, \texttt{"failed"}).
    
    \item \textbf{Register Worker:} A worker registers with the scheduler, providing hardware specifications (CPU cores, memory size, disk size, CPU frequency, OS). The system validates the worker token and creates a worker record.
    
    \item \textbf{Poll for Job:} A registered worker requests a job assignment. The scheduler selects an appropriate job using the configured scheduling algorithm and assigns it to the worker.
    
    \item \textbf{Confirm Job Received:} After receiving a job assignment, the worker confirms receipt, updating the job status to \texttt{"scheduled"}.
    
    \item \textbf{Report Job Started:} The worker reports that job execution has begun, along with a resource snapshot (CPU and memory usage at start time).
    
    \item \textbf{Execute Job:} The worker creates an isolated environment (cgroup), forks the executor process, and monitors execution.
    
    \item \textbf{Report Worker Metrics:} The worker periodically sends system metrics (CPU, memory, disk, network, load average) to the scheduler.
    
    \item \textbf{Report eBPF Metrics:} The worker sends job-specific eBPF metrics (syscall counts, I/O bytes, network bytes, CPU usage, memory usage) at regular intervals.
    
    \item \textbf{Report Job Result:} Upon job completion, the worker reports the result (success/failure, message, artifact URL) to the scheduler.
    
    \item \textbf{Manage Tokens:} The administrator generates, lists, and revokes authentication tokens for both clients and workers.
    
    \item \textbf{Query Job Timeseries:} An authorized client or administrator retrieves historical eBPF metrics for a completed job.
\end{enumerate}

\subsubsection{Non-Functional Requirements}

The following non-functional requirements define the quality attributes of the system:

\begin{description}
    \item[Scalability:] The system should support adding new worker nodes without requiring modifications to the scheduler or existing workers. The database backend (PostgreSQL) should handle increasing volumes of job and metrics data.
    
    \item[Reliability:] The system must handle worker crashes gracefully. Workers should recover incomplete jobs on restart and report them as failed. The scheduler should continue operating even if individual workers become unavailable.
    
    \item[Security:] All inter-component communication must be authenticated using bearer tokens. Tokens must be revocable, and different token types (\texttt{client} vs. \texttt{worker}) must enforce different access permissions.
    
    \item[Performance:] The scheduler should be capable of handling concurrent job submissions and worker requests. Job selection and assignment should complete within milliseconds.
    
    \item[Maintainability:] The system should use well-defined interfaces and design patterns to facilitate maintenance and extension. New scheduling algorithms and job types should be addable without modifying existing code.
    
    \item[Portability:] The scheduler and CLI components should be buildable and runnable on any system with a C++20 compiler and the required dependencies. Only the worker component is Linux-specific due to eBPF and cgroup requirements.
\end{description}

\subsection{Feasibility Analysis}

\subsubsection{Technical Feasibility}

The project relies on mature, well-supported technologies. C++20 is widely supported by modern compilers (GCC, Clang). gRPC and Protocol Buffers are production-grade frameworks with active community support. SQLite and PostgreSQL are battle-tested database systems. libbpf provides stable userspace APIs for eBPF program management. cgroups v2 is the default cgroup hierarchy in modern Linux distributions.

The key technical challenge is eBPF program compatibility, as BPF features depend on kernel version. The system requires Linux kernel version 4.18 or later (preferably 5.x for full cgroup v2 and BPF features). This requirement is reasonable for modern Linux deployments.

\subsubsection{Operational Feasibility}

The system is designed for ease of deployment. The scheduler and worker binaries are standalone executables with minimal configuration (scheduler address, port, authentication token). Database files are created automatically on first run. The PostgreSQL backend, if used, is configured through environment variables following standard conventions.

System administrators with basic Linux knowledge can deploy and operate the system. The CLI and Admin tools provide straightforward command-line interfaces with help text for all operations.

\subsubsection{Economic Feasibility}

All technologies used in the project are open-source and free of licensing costs:
\begin{itemize}
    \item C++20 compiler (Clang/GCC) --- open-source
    \item CMake --- open-source
    \item gRPC and Protocol Buffers --- Apache 2.0 license
    \item SQLite --- public domain
    \item PostgreSQL --- PostgreSQL license (permissive open-source)
    \item libbpf --- BSD license
    \item Linux kernel --- GPL v2
\end{itemize}

The hardware requirements are modest: any modern x86-64 machine with a Linux distribution can serve as a scheduler or worker node.

\subsubsection{Schedule Feasibility}

The project was completed within the academic timeline using an iterative development approach. The core components were developed in dependency order: shared library, protocol definitions, scheduler, worker, executors, eBPF monitor, CLI, and admin tool. Each component was tested incrementally, reducing integration risk.

\subsection{Object-Oriented Analysis}

This section presents the object-oriented analysis of the system using UML diagrams. The analysis identifies the key classes, their relationships, and their dynamic interactions.

\subsubsection{Class Diagram}

The class diagram (Figure~\ref{fig:class_diagram_analysis}) shows the main classes and their relationships in the Distributed Job Scheduler system.

\placeholderfig{class_diagram_analysis}{Analysis-level class diagram showing key classes and relationships}

The principal classes identified during analysis are:

\begin{description}
    \item[SchedulerServiceImpl:] The central gRPC service implementation. It handles all incoming RPCs and coordinates job scheduling, worker management, and metrics collection.
    
    \item[JobSelector (abstract):] Defines the interface for job selection strategies. Concrete implementations include FCFSSelector, SJFSelector, and AdaptiveSelector. This follows the Strategy design pattern.
    
    \item[ISchedulerDatabase (interface):] Provides a narrow interface for database queries needed by job selectors, following the Interface Segregation Principle.
    
    \item[SchedulerDatabase:] Implements the full database operations for the scheduler, including job CRUD, worker management, metrics storage, and analytics. Inherits from both the shared Database base class and ISchedulerDatabase.
    
    \item[WorkerClient:] The gRPC client used by worker nodes to communicate with the scheduler. Handles registration, job polling, metrics reporting, and result submission.
    
    \item[JobOrchestrator:] Manages the lifecycle of a job on the worker side, including cgroup creation, process forking, eBPF monitor management, and result collection.
    
    \item[JobExecutor (abstract):] Defines the interface for different stress-test executor types. Concrete implementations include StressCpuExecutor, StressMemExecutor, StressIOExecutor, and MixedLoadExecutor.
    
    \item[MetricsCollector:] Collects system-level metrics (CPU, memory, disk, network) by reading \texttt{/proc} filesystem entries.
    
    \item[EbpfMonitor:] Manages the lifecycle of the eBPF monitoring program, including loading the BPF object, attaching tracepoints, and polling BPF maps for metrics.
    
    \item[EbpfMetricStore:] Thread-safe accumulator for eBPF metrics snapshots, providing buffering between the eBPF reader threads and the metrics reporting thread.
    
    \item[ExecutorRegistry and JobSelectorRegistry:] Factory classes that manage the registration and creation of executor and selector implementations, respectively.
\end{description}

\subsubsection{Sequence Diagrams}

The sequence diagrams capture the dynamic interactions between objects during key system operations.

\placeholderfig{sequence_job_lifecycle}{Sequence diagram showing the complete job lifecycle from submission to completion}

\textbf{Job Lifecycle (Figure~\ref{fig:sequence_job_lifecycle}):}
\begin{enumerate}
    \item The Client submits a job via the CLI, which sends a \texttt{SubmitJob} gRPC request to the Scheduler. The Scheduler inserts a job record with status \texttt{"not started"}.
    \item The Worker periodically sends \texttt{GetJob} requests. The Scheduler's \texttt{JobSelector} selects the most appropriate job and assigns it, updating the status to \texttt{"ongoing"}.
    \item The Worker confirms receipt (\texttt{ConfirmJobReceived}), updating status to \texttt{"scheduled"}.
    \item The Worker starts execution (\texttt{ReportJobStarted}), recording a resource snapshot.
    \item The Worker forks the \texttt{JobOrchestrator}, which creates a cgroup and spawns the Executor process.
    \item The \texttt{MetricsCollector} periodically gathers system metrics and sends \texttt{ReportWorkerMetrics} to the Scheduler.
    \item The \texttt{EbpfMonitor} reads BPF maps and the \texttt{EbpfMetricStore} buffers snapshots, which are periodically sent via \texttt{ReportJobEbpfMetrics}.
    \item Upon job completion, the \texttt{JobOrchestrator} parses the executor output and invokes the \texttt{on\_completed} callback, which sends \texttt{ReportJobResult} to the Scheduler.
    \item The Scheduler saves the result and computes runtime analytics.
\end{enumerate}

\placeholderfig{sequence_worker_registration}{Sequence diagram showing the worker registration process}

\textbf{Worker Registration (Figure~\ref{fig:sequence_worker_registration}):}
\begin{enumerate}
    \item The Worker sends \texttt{RegisterWorker} with its hardware specs and authentication token.
    \item The Scheduler validates the token, creates a worker record, and returns a worker ID.
    \item The Worker caches the worker ID and token locally for crash recovery.
\end{enumerate}

\placeholderfig{sequence_metrics_reporting}{Sequence diagram showing the periodic metrics reporting flow}

\subsubsection{State Diagram}

The state diagram (Figure~\ref{fig:state_diagram}) captures the states and transitions of a job throughout its lifecycle.

\placeholderfig{state_diagram}{State diagram showing job states and transitions}

A job passes through the following states:
\begin{description}
    \item[not started:] Initial state after submission. The job is in the queue awaiting assignment.
    \item[ongoing:] A worker has been assigned to the job, but the worker has not yet confirmed receipt.
    \item[scheduled:] The worker has confirmed receiving the job assignment.
    \item[started:] The worker has begun executing the job and recorded an initial resource snapshot.
    \item[completed:] The job finished successfully.
    \item[failed:] The job finished with an error or the worker crashed during execution.
\end{description}

\subsubsection{Activity Diagram}

The activity diagram (Figure~\ref{fig:activity_diagram}) models the workflow of the job scheduling and execution process.

\placeholderfig{activity_diagram}{Activity diagram showing the scheduling and execution workflow}

The workflow proceeds as follows:
\begin{enumerate}
    \item The Client submits a job, which enters the scheduler's queue.
    \item The Worker polls for a job; the Scheduler applies the selection algorithm.
    \item If a suitable job is found, the Worker prepares the execution environment (cgroup creation, process fork).
    \item The Executor runs the stress test; the eBPF monitor collects kernel-level metrics.
    \item After execution, the Worker collects results and reports them to the Scheduler.
    \item The Scheduler updates the job status and stores analytics for future scheduling decisions.
\end{enumerate}
